% critique.tex — f165v, vv. 252–293 : commentaire critique.
% Parties textuelles du commentaire de l'auteur (latex/Aristophane-Thesmo Harmattan 2023 vv 1-641.tex, l. 2389–2685),
% texte et notes conservés ; répartition, harmonisation et passages signalés : repartition.md.
\cc{255}{σύζωσον ἀνύσας --- αἶρε νῦν στρόφιον.}{La répartition de ce vers comme de tout ce passage est très changeante d'une édition à une autre.
J'ai essayé pour ma part de respecter le découpage des vers du manuscrit de Ravenne, sans tailler ni fusionner çà et là des répliques qui ne l'étaient pas originellement.
J'ai confié à Euripide les conseils d'habillement, Agathon continuant d'adopter une attitude froide : il se contente de tendre à Mnésiloque les vêtements pour ne pas que ce rustre aille puiser directement dans sa garde-robe.
La principale question qui se soit posée est celle de la première partie du vers 255 : σύζωσον ἀνύσας --- αἶρε νῦν στρόφιον.
Il comporte un \italique{dicôlon} entre ἀνύσας et αἶρε sans pour autant indiquer de changement de locuteur... Ce \italique{dicôlon} est-il là pour signifier une simple pause le temps que Mnésiloque passe la crocote ? ou y a-t-il véritablement changement de locuteur ?
`σύζωσον ἀνύσας' me semble être une réplique au ton sec, marquant la pointe d'agacement qu'Agathon aurait pu sentir à la remarque de Mnésiloque du vers 254. Aussi lui ai-je accordé cette réplique ; puis Euripide de demander le \italique{strophion} (αἶρε νῦν στρόφιον) et Agathon de le tendre à Mnésiloque (ἰδού).}
\cc{263}{ἆρ᾽ ἁρμόσει μοι Ἀγ. χαλαρὰ γ᾽ο\auDessus{ῦ}{ν} χαίρεις φορῶν.}{Seulement, Brunck eut l'idée saugrenue de confier l'intégralité du vers 263 à Mnésiloque alors que le \italique{Ravennas} indiquait clairement le passage de témoin à Agathon pour la seconde partie, χαλαρὰ γο\auDessus{ῦ}{ν} χαίρεις φορῶν. Ainsi, c'est Mnésiloque qui prononcerait « tu aimes donc à te chausser large »... allusion aux mœurs sodomitiques d'Agathon. En voilà bien assez pour que tous les éditeurs postérieurs suivent en meute Brunck.
Le scholiaste était le premier à avoir ouvert le bal des affreux, qui glosait : <χαλαρὰ·>χαῦνα. διαβάλλει δὲ πάλιν τὸν Ἀγάθωνα ὡς χαῦνον.
Personne ne s'est avisé que si Mnésiloque avait effectivement lancé le trait à sous-entendu obscène qu'on lui impute, alors jamais Agathon ne lui aurait répondu avec bienveillance `σὺ τοῦτο γίγνωσκ᾽', « à toi de t'en aviser ».
Mais cette engeance de fouines à ordures perd tout sens critique dès qu'elle se met à flairer du graveleux.
Willems l'avait noté : Mnésiloque se trouve désappointé devant l'étroitesse des élégants souliers d'Agathon ; cette paire lui ira-t-elle : ἆρ᾽ ἁρμόσει μοι; se demande-t-il.
À ces mots, Agathon ne se sent pas de joie : ce rustre en préférerait sans doute une bien large ; c'est bien Agathon qui prononce `χαλαρὰ γοῦν χαίροις φορῶν'.
Et c'est à ce moment seulement que Mnésiloque saisit la perche ; il taquine une nouvelle fois Agathon sur ses mœurs supposées :\citComm{Μνησίλοχος --- σὺ τοῦτο γίγνωσκ᾽.\\
Mnésiloque --- Ah ça ! tu en sais quelque chose...}C'en est trop pour la patience d'Agathon : excédé, il décide de se retirer superbement avec l'eccyclème (exactement comme autrefois l'Euripide des \italique{Acharniens}).}
\cc{273}{τί μᾶλλον ἢ τὴν Ὑποκράτους ξυνοικίαν;}{Je conserve la leçon du manuscrit de Ravenne : \grec{Ὑποκράτους}, sans reprendre la correction communément admise \grec{Ἱπποκράτους}. En effet, le préfixe \grec{ὑ}- sert à construire des composés sur le mot \italique{porc} ; ainsi, pour se borner aux textes cités par Willems, on recense \grec{ὑωδία} (littéralement, la \italique{ressemblance à un porc}) et \grec{ὑηνία} (littéralement \italique{nature de porc}). 
En prononçant : \grec{τὴν ὑ-ποκράτου-ς-(σ)υ-ῳκίαν}, Mnésiloque insiste doublement sur la porcinité de l'hippocratique maison.
Voilà aussi pourquoi les enfants qui l'habitaient sont plus souvent qualifiés d'« \grec{υ-ἱοί Ἱπποκράτους} » que de « παῖδες Ἱπποκράτους ».}
\cc{289}{καὶ τὴν θυγατέρα χοῖρον ἀνδρός μοι τυχεῖν.}{Les vers 288 à 291 sont corrompus en plusieurs endroits. Occasion trop belle pour les philologues de nous infliger leur saleté morale. Et ils s'en sont donné à cœur joie : Mnésiloque userait d'obscénités pour parler de sa fille (τὴν θυγατέρα χοῖρον ; χοῖρον, \italique{le porc}, allusion selon eux au sexe de la femme), quand il ne lui souhaiterait tout bonnement pas les mœurs les plus infamantes. Et ce au beau milieu d’une prière aux deux déesses parmi les plus révérées !.. Juste courroux d'Alphonse Willems\nc{Tome II de son édition posthume d'Aristophane (1919), pp. 548-549. Cette note sur ce passage des \italique{Thesmophories}, la \nieme{8} qu'il écrivit, est reproduite en annexe.}. (Allez donc chercher dans tout Aristophane la moindre irrévérence à Déméter ou Coré.) On ne le répétera jamais assez : « Mais pour qui prennent-ils Aristophane ? » 
Las ! on ne détruit que ce que l'on remplace et le coup d'épée de Willems – qu'il avait pourtant aiguisée – porta dans l'eau car sa note était essentiellement aporétique. Prions les déesses thesmophores de conférer à la mienne assez de force pour dissiper à tout jamais ces nuées. 
Le nom ou surnom de la fille de Mnésiloque, †χοῖρον†, est évidemment corrompu.
La preuve en est que le scholiaste avait une autre leçon devant les yeux : χοιρίον\nc{La scholie : ὡς τῆς θυγατρὸς αὐτοῦ κατ' ἐπίκλησιν οὕτω καλουμένης. οἷον βοίδιον ἢ χρυσίον ἢ μυρτ<ί>ον. ἴσως δὲ τὸ αἰδοῖον τὸ γυναικεῖον αἰνίττεται.}. 
N'allons pas inventer un sobriquet obscène ; reportons-nous plutôt aux quelques éléments biographiques qui nous sont parvenus\nc{Voir le « Dossier Mnésiloque » de l'Introduction.}. Mnésiloque est le beau-père d'Euripide, et la fille qu'il lui a cédée en mariage s'appelle justement Χοιρίλη\nc{Certains manuscrits portent Χοιρίνη.} ; la correction est toute indiquée : la présence d'un nom propre (rare) a naturellement tendance à induire des fautes de transmission du texte ; de plus, la coquille consistant à substituer un Ν à un Λ est monnaie courante.}
\cc{291}{καὶ πρὸς θάληκον νοῦν ἔχειν μοι καὶ φρένας.}{« Pour ce qui est du †θάληκον†, qu'elle ait esprit et intelligence. » Le vers 291 a fait l'objet de nombreux débats ; et pour cause : le mot θάληκον est inconnu au bataillon des lexiques.
Une scholie en marge du manuscrit de Ravenne se référant probablement au \nieme{291} vers indique :\citComm{τὸν παιδαρίσκον. ἴσως δὲ παρὰ τὴν πόσθην αὐτὸ συνέθηκεν,\\
le jeune garçon. Peut-être à rapprocher du pénis (πόσθη).}Les commentateurs s'en sont remis à cette scholie pour en induire que πρὸς θάληκον était une transmission fautive d'un surnom de garçon, forgé sur le mot πόσθη, \italique{le membre}. La reconstruction encore la plus en vogue est celle de Dindorf : πρὸς θάληκον $ \to $ Ποσθαλίσκον, où Ποσθαλίσκον désignerait le fils (fictif) de Mnésiloque. 
À ce ποσθαλίσκον, Willems a opposé deux raisons\nc{Note 8, reproduite à la fin.}, que je crois fortes et fondées : « il est plus que douteux que ce dernier mot [Ποσθαλίσκον] soit grec, les diminutifs se faisant en ισκος et non en αλισκος. (On cite, il est vrai, κωραλίσκος, παιδαρίσκος et μειρακίσκος. Ce serait concluant s'il n'y avait pas l'intermédiaire de κωράλιον, παιδάριον et μειράκιον\nc{Chantraine, dans sa \italique{Formation des noms en grec ancien} aux pages 406 et suivantes, parle longuement du suffixe -ισκος sans même évoquer un possible -αλισκος ou -αρισκος ; Chantraine cite παιδαρίσκος et μειρακίσκος en expliquant que le Grec recourait à ce dernier diminutif parce que l'autre, μειράκιον, était devenu d'usage banal.}. Si des diminutifs comme ποσθάλιον et ποσθαλίσκος avaient existé, ils n'auraient pu manquer d'être relevés par d'anciens lexicographes.) »
À l'école de Willems, j'ai aussi appris à me méfier du scholiaste\nc{Je parle toujours, par abus de langage et conformément à un usage trop bien établi, \italique{du} scholiaste. Il est bien évident qu'en réalité les commentaires en marge des manuscrits sont de différents auteurs, et ont été compulsés les uns à la suite des autres postérieurement. Certaines scholies font preuve d'une grande érudition et apportent parfois des leçons meilleures que celle du texte transmis, d'autres, par contre, – la plupart, hélas ! – ne sont, au mieux, d'aucun intérêt voire souvent sources d'erreurs pour qui y accorderait trop d'importance. Pour finir de gâter l'affaire, les \italique{Thesmophories} est la pièce d'Aristophane la moins pourvue en scholies.}, cet « ignare byzantin », et faire bien davantage confiance à l'« alexandrin » Hésychius. Autres lieux, autres temps et autres sciences !
Comme l'a déjà fait remarquer Alphonse Willems, Hésychius n'a jamais compulsé de ποσθαλίσκον ou autre composé de πόσθη pouvant expliquer notre πρὸς θάληκον. 
En revanche, l'on trouve chez Hésychius à l'adresse θ 50 :
\citComm{θάλλικα· σάκκου εἶδος\\
Thallika : sorte de sac.}Voilà qui est beaucoup plus proche du texte transmis dans le \italique{Ravennas}, θάληκον. 
Pour le sens, j'ai reçu l'aide précieuse de Jean-Victor Vernhes qui « estime qu'il s'agit du sac où cette femme met des bijoux, la clef de sa maison, ses emplettes, sa carte de crédit (et bientôt, vous allez voir, celle de son mari), et qu'il symbolise d'une façon générale, ses affaires, ses intérêts, au profit desquels elle exploite son nigaud de conjoint. Elle finira par tout fourrer dans son sac ! »\nc{La glose d'Hésychius sur θαλικοποιοί laisse à penser que le sac en question était fait de cuir et qu'il ressemblait à un petit coffret -- une \italique{cassette}.}
Néanmoins, le mot θάλλικα (de θάλλιξ ?) est rare et circonscrit au lexique d'Hésychius. Il semble être une métathèse d'un mot beaucoup plus courant et qu'on trouve chez Aristophane : θυλάκιον. 
Si J.-V. Vernhes avait dû condenser son propos en grec, il aurait écrit quelque chose comme : σακίον, ἐν οἷσπερ τἀργύριον ταμιεύεται. Hé bien, c'est exactement la glose qu'on retrouve dans un passage de l'\italique{Onomasticon} de Pollux citant... les \italique{Thesmophories} d'Aristophane.
Pollux (X, 152) donne une liste de divers synonymes pour désigner une \italique{bourse}. Voici ce qu'il dit :\citComm{καὶ μιξάμενοι δ’ ἂν εἴποιμεν σκεύη, βαλάντια καὶ βαλαντίδια ὡς ἐν Αἰξὶν Εὔπολις. καὶ \textbf{θυλάκιον} δὲ καὶ θυλακίσκον·\\
Ἀριστοφάνης γοῦν ἐν Τριφάλητι τοῦτο ὑποδηλοῖ, ὅταν φῇ
\citComm{ἔπειτ' ἐπὶ τοὖψον ἧκε, τὴν σπυρίδα λαβών\\
καὶ θυλακίσκον καὶ τὸ μέγα βαλάντιον. }
$[$ ... $]$
ἐπὶ τοῦ αὐτοῦ δὲ καὶ σακίον, ὅταν φῇ ἐν Θεσμοφοριαζούσαις 
\citComm{\textbf{σακίον, ἐν οἷσπερ τἀργύριον ταμιεύεται.}}
Et les mêlant nous dirions, comme dans les \italique{Chèvres} d'Eupolis, les ustensiles, les grandes bourses et les petites bourses. Ainsi que θυλάκιον et θυλακίσκον. Aristophane l'illustre dans \italique{Triphalès}, lorsqu'il dit :
\citComm{Ensuite il arriva aux hors-d'œuvre, tenant la corbeille, un petit sac (θυλακίσκον) et une grande bourse.}
Ce mot désigne aussi une petite bourse, comme cela apparaît dans les \italique{Thesmophories} :
\citComm{\italique{une petite bourse, dans lesquelles on dépose l'argent.}}}Aristophane a écrit deux pièces intitulées \italique{Les Thesmophories}. Ce passage a été considéré comme un fragment des \italique{Thesmophories} perdues (fr. 343 de l'édition Kassel-Austin). En fait, ce petit bout de phrase explicatif, σακίον· ἐν οἷσπερ τἀργύριον ταμιεύεται, ressemble bien davantage à une glose d'un commentateur. Il arrive fréquemment à Pollux de fusionner des gloses avec le texte d'Aristophane lui-même\nc{Voir au commentaire du vers 1175 un autre exemple.}. Ce n'est pas en réalité aux \italique{Thesmophories} perdues mais bien au vers 291 des \italique{Thesmophories} conservées que se rattache cette glose, en tant que définition de θυλάκιον :\citComm{$<$θυλάκιον·$>$ σακίον, ἐν οἷσπερ τἀργύριον ταμιεύεται.}Une métathèse ancienne aura donné la leçon fautive du \italique{Ravennas}, comportant †θάληκον† en place de θυλάκιον, tandis que cette métathèse adjointe à sa glose `σακίον' rejoignait le lexique d'Hésychius sous la forme : θάλλικα· σάκκου εἶδος. 
Pour séduisante que soit cette reconstruction, elle n'en demeure pas moins conjecturale ; le \italique{Ravennas} ne comportant pour les \italique{Thesmophories} que très peu de scholies, et parmi les moins intéressantes.
Finissons d'ailleurs cette note en montrant que ma conjecture est aussi compatible avec l'unique scholie conservée, τὸν παιδαρίσκον. ἴσως δὲ παρὰ τὴν πόσθην αὐτὸ συνέθηκεν. Le scholiaste ayant devant lui θάληκον -- qu'il lisait indifféremment θάλικον -- devait l'analyser comme un mot formé sur τὸ θάλος, \italique{la jeune pousse} et, dans un sens métaphorique, l'\italique{enfant}, suivi d'un suffixe en -ικος ; l'on comprend le quiproquo, facilité par l'évocation deux vers plus haut de la fille de Mnésiloque et d'un imaginaire sexe féminin (χοιρίον).}
